\chapter{Аналитическая часть}

\section{Постановка задачи}
Даны матрицы $A$ размера $M \times N$ и $B$ размера $N \times Q$, содержащие целочисленные элементы. Требуется вычислить результирующую матрицу $C = A \times B$ размера $M \times Q$, элементы которой определяются соотношением:
\begin{equation}
    C_{i,j} = \sum_{k=0}^{N-1} A_{i,k} \cdot B_{k,j}, \quad i = 0, \ldots, M - 1, \quad j = 0, \ldots, Q - 1.
\end{equation}
Операция умножения матриц определена тогда и только тогда, когда количество столбцов матрицы $A$ совпадает с количеством строк матрицы $B$.

Для решения поставленной задачи требуется разработать и программно реализовать следующие алгоритмы:
\begin{itemize}
    \item стандартный алгоритм умножения матриц;
    \item базовый алгоритм Винограда;
    \item оптимизированный по варианту алгоритм Винограда.
\end{itemize}

Требуется сравненить разработанные реализацие между собой по теоретической трудоёмкости и по данным замеров процессорного времени их выполнения на квадратных матрицах ($M = N = Q$), с раздельным рассмотрением для чётных и нечётных линейных размеров матриц.

Входные данные: линейные размеры матриц $M, N, Q \in \mathbb{N}$ (для квадратных матриц --- линейный размер $N$) и целочисленные элементы матриц $A_{i,k}, B_{k,j} \in \mathbb{Z}$. Выходные данные: матрица $C$ размера $M \times Q$ ($N \times N$ для квадратных матриц), элементы которой $C_{i,j} \in \mathbb{Z}$.

\section{Вывод}
В результате аналитической части была формализована задача умножения матриц, определены форматы входных и выходных данных, а также сформирован перечень алгоритмов, подлежащих разработке, реализации и сравнительному анализу.